%% appE-further-reading \dash{} Where to Go Next
%
% NOTHING HERE IS WRITTEN FROM MEMORY, and that constraint shaped the appendix
% more than anything else. Sections E.1 and E.2 are the competitive survey in
% notes/02-grounding-and-traps.md section 1, which was run before the book was
% planned. It sources its claims about the world; its Sources section names a
% review for one of the two quotations carried here and, for Murphy's
% "encyclopedic reference", names none, so the appendix does not promise one.
% Section E.3 is notes/01-curriculum.md section 13.
%
% Section E.4 carries NAMES, and a name there is a verdict only where E.1 has
% already given one -- Strang, Deisenroth/Faisal/Ong, Goodfellow/Bengio/
% Courville, Murphy and Bishop all recur. THAT OVERLAP IS THE POINT OF THE
% WARNING BOX'S WORDING: it used to say flatly that the list below was not read
% and compared, which its own names falsified. State the rule, never a tally.
% For the rest the repository records a title and nothing more, so writing a
% characterisation would be inventing one -- the class this book spends
% forty-seven programs telling the reader not to accept.
%
% Every claim this appendix makes about a PROGRAM was checked against the
% written program before it was imported, which is the discipline Appendix D's
% pass earned. That is not gated by anything; nothing here can be.
%
%
% SECTION E.5 IS THE ROUTE INDEX FOR THE RIGOUR BOXES, and it is gated:
% tools/check_structure.py --rigour fails when a box defers outside the book
% and no row here answers it. A box that defers INSIDE the book cites the
% program it defers to, so a box with no \ref{prog:...} is outward-facing by
% construction -- that criterion is exact, it is edition-stable, and it is
% what the check reads. What it CANNOT read is whether the destination named
% contains the proof; five rows say "not surveyed" precisely because nobody
% has done that reading.
%
% STILL OUTSTANDING, AND NARROWER THAN IT WAS: there is still no \ref{app:E}
% anywhere in programs/, so a reader standing at a rigour box has no way to
% know section E.5 exists. Closing that means editing every deferring program
% and is a pass of its own. What is no longer true is that NOTHING references
% the appendix: the claims box in frontmatter/*/how-to-use.tex now points at
% section E.6, which is the only route into this appendix the book has.
%
% SECTION E.6 IS A LIST AND NOT A TABLE, deliberately, and it must not be
% tidied into one. A center+tabularx block cannot break across a page, and
% this appendix has already produced the largest overfull vbox in the book
% from exactly that -- 284 pt, from section E.3's exclusions table. What
% decides it is rows times row height, and a citation is three typeset lines
% where E.5's rows are one, so six of them come to about the height that
% overflowed. reslist is an itemize, so it breaks; the class is gone rather
% than made less likely, on both installations and at any page position.
%
% It is the book's own reference-matter list -- Appendix C replays every
% \result{} through it -- and its ragged-right rationale, written in the
% preamble beside it, is about reference matter generally. Reusing it couples
% this section to that one: restyle reslist for Appendix C and this changes
% too, which is worth knowing before anybody does.
%
% AND THE SCOPE STATEMENT HAS TWO SOURCES, not one. Sections E.3's first two
% tables are notes/01-curriculum.md section 13; the third table and the four
% "least sure about" paragraphs are notes/02-grounding-and-traps.md section 4,
% which is a second and larger exclusions list that the first draft of this
% appendix never opened. If you extend the scope statement, open both.
%
% NO WEB ADDRESSES. The body of this book prints none, and a printed address
% rots faster than a title does.

\chapter{Where to Go Next}
\label{app:E}

\noindent
This appendix is in two halves, and the difference between them is stated
rather than hidden.

Sections~\ref{sec:E-survey} and~\ref{sec:E-positions} carry a verdict on every
entry, because those works were read and compared before this book was
planned: the survey behind them exists to answer two questions \dash{}
\emph{what already exists}, and \emph{what does each do well and badly for this
reader?} Where a quotation here is not attributed to a named review, that is
because the record behind it does not name one.
Section~\ref{sec:E-not} says what this book does not teach, why not, and where
the subject is done properly.

Section~\ref{sec:E-depth} is a list of names, and it says which of them carry a
verdict. Some are works Section~\ref{sec:E-survey} has already judged, and that
judgement stands; the rest this book has not held to that standard, and a
one-line characterisation invented to fill the gap would be exactly the thing
forty-seven programs have been asking you not to accept.

No web addresses appear here. A title and an author outlive a URL.

\section{The books this one was measured against}
\label{sec:E-survey}
\index{further reading!competitive survey}

Ordered by how close each work sits to the reader this book was written for: a
working engineer, competent in code, whose school mathematics has decayed, and
who is now shipping models.

Every entry says what the work does well, what it does badly \emph{for this
reader}, the floor it assumes before its first page, and whether it is a first
reading or a reference. \textbf{None of them says how long it takes}, and the
omission is deliberate rather than an oversight: nobody on this project has
timed a reading, and a duration invented to fill the column would be the first
unsupported claim in the book. The floor and the scale are what decide it, and
those are facts.

\medskip
\textbf{Deisenroth, Faisal and Ong, \emph{Mathematics for Machine Learning}.}
The closest competitor, and it names the same gap this book does \dash{} its
authors write that the distance between school mathematics and the level needed
to read a machine-learning textbook is too big for many people. \textbf{Does well:}
the single best-organised map of the territory, free as a PDF, with a coherent
narrative from vector spaces to four worked methods, and the best short
treatment of matrix decompositions available. \textbf{Does badly, for this
reader:} it states the school-to-machine-learning gap as its target and then
starts above it \dash{} chapter 2 opens on systems of linear equations, and
nothing in it teaches you to rearrange an inequality or shift a summation
index \dash{} the floor it assumes in practice is fluent algebra and some
calculus. Its four target methods are pre-deep-learning, there is not one page
on overflow or conditioning, and most exercises have no worked solution.
\textbf{Verdict:} the best reference to point at, and not a path from zero
despite its preface.

\medskip
\textbf{Goodfellow, Bengio and Courville, \emph{Deep Learning}, Part~I.}
\textbf{Does well:} chapter 4, \emph{Numerical Computation}, is the most
valuable chapter in this whole survey and the one nobody imitates
\dash{} overflow and underflow, the log-sum-exp stabilisation, conditioning,
the Jacobian and the Hessian. It is the only work here that treats the machine
as a finite-precision object. \textbf{Does badly:} Part~I is a refresher rather
than a teacher, four chapters standing in for two years of undergraduate
mathematics, which is a refresher for somebody who has met that material and
not a teacher for somebody who has not, with no exercises and no worked
examples; and it is frozen in
2016, so there is no transformer and no modern mixed-precision practice.
\textbf{Verdict:} that chapter is the ancestor of this book's Part~II, and this
book goes an order of magnitude deeper with numbers.

\medskip
\textbf{Strang, \emph{Linear Algebra and Learning from Data}.}
\textbf{Does well:} Strang is the finest expositor of linear algebra alive, and
this is where he connects the four subspaces, the SVD and low-rank
approximation to data. Putting the SVD at the centre of gravity is exactly
right. \textbf{Does badly:}
reviewers are blunt about the density \dash{} the MAA review says he has tried
to cram far too much material into 432 pages, so each topic gets a very quick
presentation and results are stated without proof or citation. It also assumes
a completed first course; it is a second one.
\textbf{Verdict:} a source of framing rather than a model of construction.

\medskip
\textbf{Boyd and Vandenberghe, \emph{Introduction to Applied Linear Algebra}.}
\textbf{Does well:} the best-designed applied linear algebra text in this
survey, and the closest to this book in temperament. It is aggressively narrow
\dash{}
built on linear independence and the QR factorisation and almost nothing else
\dash{} and it proves that a book can omit most of its subject and be better
for it. Free PDF, with Python and Julia companions.
\textbf{Does badly, for this reader:} by design it omits eigenvalues, the SVD
and determinants, which are precisely the parts needed for conditioning,
attention and covariance, and it still assumes you can do the algebra.
\textbf{Verdict:} the scoping discipline to imitate. Its omissions are not this
book's omissions.

\medskip
\textbf{Boyd and Vandenberghe, \emph{Convex Optimization}.}
\textbf{Does well:} the reference. Convexity, duality, and above all the
recognition problem \dash{} \emph{is my problem convex?} \dash{} are
permanently useful vocabulary. Free PDF, extensive exercises.
\textbf{Does badly, for this reader:} it wants good linear algebra up front by
its own statement, and deep learning is not in its class of problems.
Non-convex, stochastic, over-parameterised objectives sit outside the theory,
and its convergence-rate proofs do not transfer.
\textbf{Verdict:} take the vocabulary into Programs~\ref{prog:P19}
and~\ref{prog:P22}; do not take the curriculum.

\medskip
\textbf{Bishop, \emph{Pattern Recognition and Machine Learning}, and Bishop and
Bishop, \emph{Deep Learning: Foundations and Concepts}.}
\textbf{Does well:} the organisation, the figures and the probabilistic-model
material are among the best written anywhere, and the 2024 volume brings the
same treatment to transformers and diffusion. \textbf{Does badly:} both are
more advanced than most other introductions and assume multivariable calculus,
linear algebra and probability as working tools; the maths is always in service
of a model and is never taught from the floor.
\textbf{Verdict:} a destination rather than a route \dash{} and this book's own
success condition is that you can now open it.

\medskip
\textbf{Murphy, \emph{Probabilistic Machine Learning} \dash{} the 2022
introduction and the 2023 advanced volume.}
\textbf{Does well:} the most complete and most current single treatment, unified
by probabilistic modelling, with a good mathematical-background section and open
code. \textbf{Does badly:} graduate-level by its own statement, assuming
calculus, statistics and linear algebra; a reviewer's description of it as an
encyclopedic reference is both praise and the problem, at over a thousand pages
across the two.
\textbf{Verdict:} the reference to route to per topic, once you can read it.

\medskip
\textbf{Blum, Hopcroft and Kannan, \emph{Foundations of Data Science}.}
\textbf{Does well:} the only book in this list that takes seriously the thing
engineers most reliably get wrong \dash{} how data behaves in high dimensions.
Concentration of measure, random projection, the SVD, random walks.
\textbf{Does badly, for this reader:} it is a theory course, proof-driven and
aimed at computer-science theory students, and its progression does not
resemble a practitioner's need.
\textbf{Verdict:} the source for Program~\ref{prog:P05}'s content, and this book
translates it out of theorem-and-proof form into something you can act on.

\section{Two positions, and one method}
\label{sec:E-positions}
\index{further reading!positions}

\medskip
\textbf{3Blue1Brown.} \emph{Essence of Linear Algebra} and \emph{Essence of
Calculus} do something the books here do not: they make the geometric content
of a determinant or a derivative visible in seconds.

And the strongest argument for reading a book like this one instead is made by
their author. Sanderson has said publicly that he worries his work is
intellectual candy, offering the feeling of understanding, and that unless the
viewer actively engages \dash{} reinvents it, does problems \dash{} it is not
going to stick; he extends the point to Feynman's lectures, saying that how
satisfying content is to consume does not correlate with long-term learning.

That is a concession by the person with the most to lose from it, and it is
this book's own thesis in one sentence. The right relationship is complementary:
watch the video for the picture, work the program for the retention.

\medskip
\textbf{fast.ai, and the position that you need less mathematics than you
think.} It is right about \emph{entry}, and this book concedes the point
without reservation: the claim that two years of mathematics must come before
anyone may train a classifier kept people out of the field, produced nothing,
and is empirically false.

What the position answers is \emph{can I ship something?} What it does not
answer is why the loss went to \code{nan} at step four thousand and only in
half precision; whether a three-point evaluation gain is real or one draw from
a noisy distribution; or whether a gradient-accumulation setting quietly changed
the objective. Those questions arrive \emph{after} you ship, and they are
mathematics questions.

So: do not learn the mathematics first. This book is what you reach for second,
when something is wrong and you cannot see why.

\medskip
\textbf{Stroud, \emph{Engineering Mathematics} \dash{} the method, not a
competitor.} The frame, the answer overleaf, the Quiz, the Summary and the
\emph{Can you?} checklist are all his. Its own limitations are this book's
obligations, and they are named in the introduction rather than hidden:
no rigour; recipe-level linear algebra with no vector spaces and no
decompositions; descriptive statistics with no inference; no discrete
mathematics; no optimisation in the first volume; and \dash{} the complaint
readers make most \dash{} useless as a reference once read. Part~III answers
the linear algebra, Part~IV the discrete mathematics and the rigour, Part~VI
the optimisation and Part~VII the inference; Appendices~B, C and D exist for
the last.

\section{What this book does not teach, and where to go instead}
\label{sec:E-not}
\index{further reading!exclusions}

The house rule is that the book may say a topic is not worth its cost to this
reader. It may not pretend the topic does not exist.

First the rigorous foundations underneath what this book teaches \dash{} the
places it borrows a result and does not prove it.

\begin{center}
\begin{tabularx}{\linewidth}{@{}XXl@{}}
\toprule
\textbf{Not taught} & \textbf{Why not} & \textbf{Where to go} \\
\midrule
Measure-theoretic probability & Nothing this book pays off needs a
  sigma-algebra, and the machinery costs a semester & Williams; Durrett \\
Real analysis and proof at the epsilon--delta level & Program~\ref{prog:P14}
  teaches you to \emph{read} a theorem and deliberately not to write one & Abbott;
  Tao \\
Partial differential equations & The audience meets them in
  physics-informed work, which a shallow treatment does not serve & Strauss;
  Evans \\
Stochastic differential equations, and so the full mathematics of diffusion
  models & The book gives the Gaussian and the fact that variances add, which
  is the whole of a discrete-time chain of Gaussians, and never applies them to
  one & Särkkä and Solin \\
Complex analysis, and the Fourier transform, which this book does not mention &
  Program~\ref{prog:F08} gives the unit circle; convolution theorems are a book
  of their own & Bracewell \\
\bottomrule
\end{tabularx}
\end{center}

\medskip
And the specialisms that begin where a working depth ends. Each of these is a
field, and the book stops at the point where an engineer stops needing to
argue about it.

\begin{center}
\begin{tabularx}{\linewidth}{@{}XXl@{}}
\toprule
\textbf{Not taught} & \textbf{Why not} & \textbf{Where to go} \\
\midrule
Numerical linear algebra at implementation level & Program~\ref{prog:P09} says
  why not to invert and Program~\ref{prog:P11} why the normal equations square
  the condition number; neither teaches you to implement a QR factorisation &
  Trefethen and Bau; Higham \\
Convex optimisation theory, duality in full & Programs~\ref{prog:P19}
  and~\ref{prog:P22} give the working subset, and there is no case for
  paraphrasing one of the best-written books in mathematics & Boyd and
  Vandenberghe \\
Statistical learning theory: VC dimension, generalisation bounds &
  Program~\ref{prog:P14} names them as claims to read carefully; the bounds are
  almost always vacuous at realistic scale & Shalev-Shwartz and Ben-David \\
Reinforcement learning theory and convergence results &
  Programs~\ref{prog:P21} and~\ref{prog:P22} give the gradient estimators and
  the KL constraint, which is what an engineer touching alignment work actually
  manipulates & Sutton and Barto \\
Category theory, differential geometry, manifolds & Occasionally invoked in
  interpretability writing; not load-bearing for the work, and no source here
  was surveyed for the category theory & Lee, for the geometry \\
\bottomrule
\end{tabularx}
\end{center}

\medskip
And the classical curriculum this book declines. The two tables above are
subjects it \emph{borrows} from, so each names a book. These are techniques a
mathematics course would drill, and their destination is not a book at all
\dash{} it is a program of this one and a routine you already have installed.
That is the whole of the argument for cutting them: not that they are
uninteresting, but that the work they were invented to do by hand is now done
by a library, and the skill that survives is knowing which library call and
why.

\begin{center}
\begin{tabularx}{\linewidth}{@{}XX@{}}
\toprule
\textbf{Not drilled} & \textbf{What does the work instead} \\
\midrule
Integration technique: by parts, by substitution, partial fractions & Nothing
  in a training loop integrates symbolically. An integral arrives as an
  expectation and is estimated by sampling, which is Program~\ref{prog:F13}
  and then Program~\ref{prog:P25} \\
Characteristic polynomials, Cayley--Hamilton, Jordan form & No library computes
  an eigenvalue from a polynomial. Program~\ref{prog:P10} does eigenvalues as
  invariant directions, and Program~\ref{prog:P11} does the decomposition you
  actually call \\
Determinants by cofactor expansion, Cramer's rule & Factorially expensive and
  never used at size. Program~\ref{prog:P09} teaches what a determinant
  \emph{means}, gives the $2 \times 2$ formula, and says that beyond it a
  machine computes one by elimination and so should you \\
Frequentist estimator theory: sufficiency, the Cramér--Rao bound, UMVU &
  Program~\ref{prog:P26} does maximum likelihood as an \emph{objective}, which
  is the part that shows up in every loss function; the optimality theory does
  not reach the desk \\
The classical test zoo: $t$ tables, $\chi^{2}$, ANOVA, $F$ & You have a
  computer. Program~\ref{prog:P27} answers the same questions by resampling and
  by an exact test on the discordant pairs, with fewer assumptions and no table
  to look up \\
\bottomrule
\end{tabularx}
\end{center}

\medskip
\textbf{And four exclusions this book is least sure about}, stated because
saying so now is cheaper than defending them later.

\emph{Convex optimisation theory} is out on the grounds that deep learning is
not convex, and the counter-argument is good: quantisation, calibration,
serving and routing problems and constrained decoding are convex, and an
engineer who cannot recognise one solves it with a loop.
Programs~\ref{prog:P19} and~\ref{prog:P22} keep the recognition skill for
exactly that reason, and this is the first exclusion to revisit if readers
report needing more.

\emph{Discrete mathematics} stops well short of algorithms and complexity.
That is defensible only because the reader is a software engineer who has met a
graph before \dash{} an assumption this book makes nowhere else.

\emph{Causal inference} is excluded outright, and that is increasingly wrong as
evaluation and A/B analysis become part of the job. It is a candidate for a
second edition rather than a subject that does not belong.

\emph{Stochastic differential equations} are excluded in the table above. That
is a defensible line today and it may not be one in a few years.

\section{The subjects it does teach, but only to working depth}
\label{sec:E-depth}
\index{further reading!working depth}

\begin{warning}
\textbf{These are names, and a name here is a verdict only where
Section~\ref{sec:E-survey} has already given one.} That section says what each
work does well and badly because those works were read and compared against
each other before this book was planned. Where a name below appears there too,
its verdict stands and this section adds nothing to it. The rest were not read
and compared, and inventing a characterisation for them is precisely the move
this book has spent forty-seven programs asking you to refuse \dash{} so for
those, what follows is a pointer and nothing more.
\end{warning}

\medskip
For \textbf{linear algebra} beyond the working depth of Part~III: Strang, and Axler
for the abstract treatment. For the same ground with a machine-learning slant,
Deisenroth, Faisal and Ong.

For \textbf{probability}, Blitzstein and Hwang. For \textbf{inference}, Wasserman,
\emph{All of Statistics}. For \textbf{Bayesian methods} past
Program~\ref{prog:P28}, Gelman and colleagues, \emph{Bayesian Data
Analysis}.

For \textbf{information theory} past Part~VIII, Cover and Thomas. For
\textbf{numerical optimisation} as a subject in its own right rather than as
the update rules Program~\ref{prog:P20} derives, Nocedal and Wright.

And for the machine learning that consumes all of it: Goodfellow, Bengio and
Courville for the 2016 picture, Murphy and Bishop for the current one.

\medskip
And one gap here is worth naming rather than leaving to be found. Part~II goes
an order of magnitude deeper than the one chapter this survey found on the
subject, and no source for floating-point arithmetic itself was surveyed for
this appendix, so this section has no pointer for it.
Section~\ref{sec:E-not} sends you to Trefethen and Bau and to Higham for
numerical linear algebra, which is the neighbouring subject rather than that
one.

\section{Where each unproved statement goes}
\label{sec:E-route}
\index{further reading!rigour boxes}
\index{rigour boxes!destinations}

A rigour box names a result, says it is not proved here, and says where the
proof lives. Most of them point at another program of this book. The ones that
point \emph{outside} it name a kind of course \dash{} \enquote{any first
analysis course} \dash{} and, with a single exception, no title. This section
is the index for them, and the rule rather than the count is what is kept: a
box that leaves the book has a row here.

\begin{note}
\textbf{A destination here is a work this appendix has already introduced, or
it is nothing, and the column says which.} Where it says \emph{not surveyed},
no source in this appendix was read for that subject, and naming one anyway
would be the move forty-seven programs have asked you to refuse. The last of
the three tables below is the rows that say it. That is the state of the
record rather than a gap in the typesetting, and it is worth knowing before
you go looking.
\end{note}

\medskip
First the foundations underneath the mathematics this book does teach.

\begin{center}
\begin{tabularx}{\linewidth}{@{}Xll@{}}
\toprule
\textbf{Stated and not proved} & \textbf{Deferred from}
  & \textbf{Where to go} \\
\midrule
When a limit exists, and when a covering does not matter &
  \ref{prog:F04}, \ref{prog:F13}, \ref{prog:P09} & Abbott; Tao \\
Differentiability, and Taylor with a remainder &
  \ref{prog:P15}, \ref{prog:P17} & Abbott; Tao \\
That two bases have the same size, and the spectral theorem &
  \ref{prog:P04}, \ref{prog:P10} & Axler; Strang \\
What a probability is on a continuum &
  \ref{prog:P23} & Williams; Durrett \\
That independent things factorise, and characteristic functions &
  \ref{prog:P25} & Blitzstein and Hwang \\
\bottomrule
\end{tabularx}
\end{center}

\medskip
Then the results the book uses and does not derive.

\begin{center}
\begin{tabularx}{\linewidth}{@{}Xll@{}}
\toprule
\textbf{Stated and not proved} & \textbf{Deferred from}
  & \textbf{Where to go} \\
\midrule
Constraint qualification, and sufficiency &
  \ref{prog:P22} & Boyd and Vandenberghe \\
Kraft's inequality, and the source-coding bound &
  \ref{prog:P29}, \ref{prog:P30} & Cover and Thomas \\
The evidence lower bound, and variational inference &
  \ref{prog:P19} & Murphy \\
Asymptotic efficiency, and the Cramér--Rao bound &
  \ref{prog:P26} & excluded above \\
What a residual path and a normalisation layer do to the variance &
  \ref{prog:P25} & Program~\ref{prog:P32} \\
\bottomrule
\end{tabularx}
\end{center}

\medskip
And the five this appendix cannot discharge.

\begin{center}
\begin{tabularx}{\linewidth}{@{}Xll@{}}
\toprule
\textbf{Stated and not proved} & \textbf{Deferred from}
  & \textbf{Where to go} \\
\midrule
Stars and bars, and counting with repetition &
  \ref{prog:P12} & not surveyed \\
How many nearly perpendicular directions a space really holds &
  \ref{prog:P05} & not surveyed \\
Why resampling estimates a sampling distribution &
  \ref{prog:P27} & not surveyed \\
How to write a proof rather than read one &
  \ref{prog:P14} & not surveyed \\
The repeated-root case of a linear recurrence &
  \ref{prog:P20} & not surveyed \\
\bottomrule
\end{tabularx}
\end{center}

\medskip
Two rows of that last table have a pointer in the program itself rather than
here: Program~\ref{prog:P27} names the paper the bootstrap came from, and
Program~\ref{prog:P05} names the field its question belongs to. The rest name
a kind of course and no title. \textbf{One of them is not an omission but the
book's own position}: this book teaches you to read a theorem and says in as
many words that you will not need to write one, so a proof-writing text is the
one destination it declines to send you to.

\medskip
That last sentence is the point of everything above it. This book teaches none
of those subjects to the depth its own sources do, and it was never meant to.
What it claims is narrower and checkable: that you can now open them.

\section{Where each attributed claim came from}
\label{sec:E-sources}
\index{further reading!sources}
\index{sources!for attributed claims}

\textbf{This section does a different job from the rest of the appendix, and
the two are kept apart for that reason.} Everything above recommends: it says
where to go next and carries a verdict on each destination. This section
attributes: it names the work behind each claim this book makes about somebody
else's result, so that a reader who wants to check one can. A work you are
being sent to read and a work a sentence rests on are not the same kind of
thing, and an entry here is neither a recommendation nor a warning against.

\begin{note}
\textbf{Volume and page range, and no issue numbers.} An issue number is one
more place for a transcription error to hide, and it buys nothing: a journal, a
volume and a page range locate an article in any index in the world. Nothing is printed here that was not checked
against the journal's own record, which is the rule the rest of this book runs
on and the reason the list is short. There are no web addresses, for the reason
this appendix already gives; where a work was never published outside a
preprint server its identifier is given instead, which is a document number
rather than an address.
\end{note}

\medskip
First the evidence for the method itself. This is the argument \emph{How to use
this book} makes for the format it puts you through, and it is the one place
where a reader wanting to check the book's case for its own existence has to
leave the book to do it.

\begin{reslist}
\item K. A. Stroud, \emph{Engineering Mathematics: Programmes and Problems},
  Macmillan, 1970 \dash{} the format itself.
\item C. C. Kulik, B. J. Schwalb and J. A. Kulik, \enquote{Programmed
  Instruction in Secondary Education: A Meta-Analysis of Evaluation Findings},
  \emph{Journal of Educational Research}, 75, 133--138, 1982 \dash{} the case
  against it.
\item H. L. Roediger and J. D. Karpicke, \enquote{Test-Enhanced Learning:
  Taking Memory Tests Improves Long-Term Retention}, \emph{Psychological
  Science}, 17, 249--255, 2006 \dash{} retrieval practice, and the delays at
  which the advantage appears.
\item N. Kornell, M. J. Hays and R. A. Bjork, \enquote{Unsuccessful Retrieval
  Attempts Enhance Subsequent Learning}, \emph{Journal of Experimental
  Psychology: Learning, Memory, and Cognition}, 35, 989--998, 2009 \dash{}
  errorful generation.
\item B. Butterfield and J. Metcalfe, \enquote{Errors Committed with High
  Confidence Are Hypercorrected}, \emph{Journal of Experimental Psychology:
  Learning, Memory, and Cognition}, 27, 1491--1494, 2001 \dash{} why the
  benefit is largest when the learner was confident.
\item J. D. Karpicke, A. C. Butler and H. L. Roediger, \enquote{Metacognitive
  Strategies in Student Learning: Do Students Practise Retrieval When They
  Study on Their Own?}, \emph{Memory}, 17, 471--479, 2009 \dash{} why the
  method that teaches more feels worse.
\end{reslist}

\medskip
Then the results the programs name. Each of these is a place where a program
says whose work a line came from, and until now said it with a surname and no
title.

\begin{reslist}
\item C. Eckart and G. Young, \enquote{The Approximation of One Matrix by
  Another of Lower Rank}, \emph{Psychometrika}, 1, 211--218, 1936 \dash{}
  Program~\ref{prog:P11}.
\item B. Efron, \enquote{Bootstrap Methods: Another Look at the Jackknife},
  \emph{The Annals of Statistics}, 7, 1--26, 1979 \dash{}
  Program~\ref{prog:P27}.
\item D. P. Kingma and J. Ba, \enquote{Adam: A Method for Stochastic
  Optimization}, \emph{International Conference on Learning Representations},
  2015 \dash{} Program~\ref{prog:F02}.
\item J. L. Ba, J. R. Kiros and G. E. Hinton, \enquote{Layer Normalization},
  arXiv:1607.06450, 2016 \dash{} Program~\ref{prog:F02}.
\item A. Vaswani et al., \enquote{Attention Is All You Need}, \emph{Advances
  in Neural Information Processing Systems}, 30, 5998--6008, 2017 \dash{}
  Programs~\ref{prog:F02} and~\ref{prog:F08}.
\item K. He, X. Zhang, S. Ren and J. Sun, \enquote{Delving Deep into
  Rectifiers: Surpassing Human-Level Performance on ImageNet Classification},
  \emph{IEEE International Conference on Computer Vision}, 1026--1034, 2015
  \dash{} Program~\ref{prog:P25}.
\end{reslist}

\medskip
\textbf{And one claim in the front matter has no source here.} That small steps
and linear sequencing turned out not to be essential, and that immediate
reinforcement turned out not to be critical, is a claim about the format's
mechanism rather than about its results, and the meta-analysis above does not
carry it. No source was surveyed for it, so none is named: putting up a paper
that might contain it would be the move forty-seven programs have asked you to
refuse. It is the same treatment the rows above marked \emph{not surveyed} get,
and for the same reason.

\medskip
An entry here is not a reading recommendation, and its presence says nothing
about whether the work is worth your evening. What it says is narrower, and it
is the only thing this appendix has ever claimed: that the sentence resting on
it was checked against it, and that you can now do the same.
